\documentclass[a4paper,landscape,8pt]{extarticle}
\usepackage[landscape,margin=0.7cm,top=0.5cm,bottom=0.7cm]{geometry}
\usepackage[utf8]{inputenc}
\usepackage[T1]{fontenc}
\usepackage[ngerman]{babel}
\usepackage{helvet}
\renewcommand{\familydefault}{\sfdefault}
\usepackage{xcolor}
\usepackage{tikz}
\usetikzlibrary{shapes, arrows.meta, positioning}
\usepackage{enumitem}
\usepackage{multicol}
\usepackage{array}
\usepackage{fancyhdr}
\usepackage{amssymb}
\usepackage{eurosym}

% Farben
\definecolor{appleblue}{HTML}{007AFF}
\definecolor{icloudblue}{HTML}{3693F3}
\definecolor{appstoreblue}{HTML}{1D82F5}
\definecolor{findmygreen}{HTML}{34C759}
\definecolor{imessagegreen}{HTML}{34C759}
\definecolor{lightgray}{HTML}{F5F5F5}
\definecolor{bordergray}{HTML}{BBBBBB}
\definecolor{warnred}{HTML}{C62828}

\pagestyle{fancy}
\fancyhf{}
\renewcommand{\headrulewidth}{0pt}
\fancyfoot[C]{\footnotesize Seite \thepage{} von 3}

\setlength{\parindent}{0pt}
\setlength{\parskip}{2pt}

% Kompakte Listen
\setlist{nosep, leftmargin=1em, itemsep=0pt, topsep=1pt}

\begin{document}

% ============================================================================
% SEITE 1: INFORMATIONSBLATT
% ============================================================================

\noindent{\Large\bfseries Das Apple-Ökosystem: Alles verbunden} \hfill
{\footnotesize\textbf{Lernziele:} Apple-ID verstehen | Dienste kennen | Synchronisation}

\vspace{1mm}
\hrule
\vspace{1mm}

% Drei-Spalten-Layout
\begin{minipage}[t]{0.32\linewidth}
\begin{center}
\colorbox{appleblue!15}{%
\begin{minipage}{0.95\linewidth}
\vspace{2mm}
\begin{center}
{\Large\bfseries\textcolor{appleblue}{1. Die Apple-ID}}\\[1pt]
{\footnotesize\itshape Dein Schlüssel zu allem}
\end{center}
\vspace{2mm}
\end{minipage}%
}
\end{center}

\vspace{1mm}

\textbf{Was ist die Apple-ID?}\\
Dein Konto bei Apple -- sieht aus wie E-Mail, ist aber viel mehr!

\textbf{Analogie: Wie ein Hotelgast}
\begin{itemize}
\item Apple-ID = deine Gästekarte
\item Öffnet: Zimmer (iCloud), Restaurant (Apps)
\item Funktioniert in allen ``Hotels'' (Geräten)
\end{itemize}

\textbf{Was ist verbunden?}
\begin{itemize}
\item iCloud (Fotos, Kontakte, Backup)
\item App Store (gekaufte Apps)
\item ``Wo ist?'' -- Geräte orten
\item iMessage, FaceTime
\end{itemize}

\vspace{2mm}
\fcolorbox{warnred}{warnred!5}{%
\begin{minipage}{0.92\linewidth}
\footnotesize
\textbf{Wichtig:} Wer Apple-ID + Passwort hat, kann alles sehen!
\end{minipage}%
}

\end{minipage}%
\hfill%
\begin{minipage}[t]{0.32\linewidth}
\begin{center}
\colorbox{icloudblue!15}{%
\begin{minipage}{0.95\linewidth}
\vspace{2mm}
\begin{center}
{\Large\bfseries\textcolor{icloudblue}{2. iCloud}}\\[1pt]
{\footnotesize\itshape Speicher in der Wolke}
\end{center}
\vspace{2mm}
\end{minipage}%
}
\end{center}

\vspace{1mm}

\textbf{Was speichert iCloud?}
\begin{itemize}
\item Fotos und Videos
\item Kontakte und Kalender
\item Notizen und Erinnerungen
\item Backups deines iPhones
\item Passwörter (Schlüsselbund)
\end{itemize}

\textbf{Wo finde ich es?}
\begin{itemize}
\item Einstellungen → [Dein Name] → iCloud
\item Im Browser: \textbf{icloud.com}
\end{itemize}

\textbf{Speicherplatz:}\\
5 GB kostenlos. Mehr: 50 GB = 0,99\euro{}/Monat

\vspace{2mm}
\fcolorbox{icloudblue}{icloudblue!5}{%
\begin{minipage}{0.92\linewidth}
\footnotesize
\textbf{Vorteil:} Ein Foto auf dem iPhone erscheint automatisch auf iPad und Mac!
\end{minipage}%
}

\end{minipage}%
\hfill%
\begin{minipage}[t]{0.32\linewidth}
\begin{center}
\colorbox{findmygreen!15}{%
\begin{minipage}{0.95\linewidth}
\vspace{2mm}
\begin{center}
{\Large\bfseries\textcolor{findmygreen}{3. Weitere Dienste}}\\[1pt]
{\footnotesize\itshape Apps, Nachrichten, Finden}
\end{center}
\vspace{2mm}
\end{minipage}%
}
\end{center}

\vspace{1mm}

\textbf{App Store:}
\begin{itemize}
\item Der ``Laden'' für Apps
\item Einmal gekauft = auf allen Geräten
\end{itemize}

\textbf{``Wo ist?'':}
\begin{itemize}
\item Zeigt Geräte auf der Karte
\item Kann klingeln lassen
\item Sperren oder löschen (Notfall!)
\end{itemize}

\textbf{iMessage \& FaceTime:}
\begin{itemize}
\item \textbf{Blaue Blasen} = iMessage (kostenlos)
\item Grüne Blasen = SMS (kann kosten)
\item FaceTime = kostenlose Videoanrufe
\end{itemize}

\vspace{2mm}
\fcolorbox{findmygreen}{findmygreen!5}{%
\begin{minipage}{0.92\linewidth}
\footnotesize
\textbf{Tipp:} iPhone verloren? Über icloud.com orten!
\end{minipage}%
}

\end{minipage}

\vspace{1mm}

% Zusammenfassung
\begin{center}
\fcolorbox{bordergray}{lightgray}{%
\begin{minipage}{0.98\linewidth}
\centering
\vspace{1.5mm}
\textbf{Zusammenfassung:}
\textcolor{appleblue}{\textbf{Apple-ID}} = Schlüssel zu allem ~$|$~
\textcolor{icloudblue}{\textbf{iCloud}} = Speicher + Synchronisation ~$|$~
\textcolor{findmygreen}{\textbf{``Wo ist?''}} = Geräte finden ~$|$~
Alles \textbf{verbunden} über alle Geräte
\vspace{1.5mm}
\end{minipage}%
}
\end{center}

\newpage

% ============================================================================
% SEITE 2: ÜBUNGEN
% ============================================================================

\begin{center}
{\LARGE\bfseries Übungen: Das Apple-Ökosystem verstehen}\\[3pt]
{\normalsize Schau dir Seite 1 an und beantworte die Fragen}
\end{center}

\vspace{3mm}
\hrule
\vspace{4mm}

\begin{multicols}{2}

\textbf{\large Teil A: Was nutze ich wofür?}

Verbinde die Situation mit dem richtigen Dienst:

\vspace{2mm}

\renewcommand{\arraystretch}{1.6}
\begin{tabular}{@{}p{7cm}p{4cm}@{}}
1. Foto vom iPhone auf dem iPad sehen: & \dotfill \\
2. Neue App herunterladen: & \dotfill \\
3. Verlorenes iPhone finden: & \dotfill \\
4. Mit Enkelin videotelefonieren: & \dotfill \\
5. Fotos sichern bei vollem Speicher: & \dotfill \\
6. Kostenlose Nachricht an Sohn: & \dotfill \\
\end{tabular}
\renewcommand{\arraystretch}{1}

{\footnotesize\textit{Lösungen: iCloud, App Store, ``Wo ist?'', FaceTime, iCloud, iMessage}}

\vspace{6mm}

\textbf{\large Teil B: Richtig oder Falsch?}

\vspace{2mm}

\renewcommand{\arraystretch}{1.6}
\begin{tabular}{@{}p{8cm}cc@{}}
& R & F \\
1. Jedes Gerät braucht eine eigene Apple-ID. & $\square$ & $\square$ \\
2. Gekaufte Apps sind auf allen Geräten nutzbar. & $\square$ & $\square$ \\
3. iCloud-Speicher ist unbegrenzt kostenlos. & $\square$ & $\square$ \\
4. Mit ``Wo ist?'' kann man das iPhone löschen. & $\square$ & $\square$ \\
5. Blaue Nachrichten sind SMS und kosten Geld. & $\square$ & $\square$ \\
\end{tabular}
\renewcommand{\arraystretch}{1}

{\footnotesize\textit{Lösungen: F, R, F, R, F}}

\columnbreak

\textbf{\large Teil C: Begriffe erklären}

Was ist was? Erkläre in eigenen Worten:

\vspace{3mm}

\textbf{Apple-ID:} \dotfill

\vspace{2mm}

\dotfill

\vspace{4mm}

\textbf{iCloud:} \dotfill

\vspace{2mm}

\dotfill

\vspace{4mm}

\textbf{Synchronisation:} \dotfill

\vspace{2mm}

\dotfill

\vspace{6mm}

\textbf{\large Teil D: Situationen}

Was machst du?

\vspace{3mm}

\textbf{Situation 1:}\\
Dein iPhone ist weg. Du hast aber einen Mac zu Hause.

\vspace{2mm}

\dotfill

\vspace{2mm}

\dotfill

\vspace{4mm}

\textbf{Situation 2:}\\
Du kaufst eine App auf dem iPhone. Ist sie auch auf dem iPad?

\vspace{2mm}

\dotfill

\end{multicols}

\newpage

% ============================================================================
% SEITE 3: CHECKLISTE
% ============================================================================

\begin{center}
{\LARGE\bfseries Checkliste: Das habe ich verstanden}\\[3pt]
{\normalsize Geh die Liste durch und hake ab, was du sicher verstanden hast}
\end{center}

\vspace{3mm}
\hrule
\vspace{6mm}

% Drei Boxen nebeneinander
\begin{minipage}[t]{0.31\linewidth}
\begin{center}
\fcolorbox{appleblue}{appleblue!10}{%
\begin{minipage}{0.92\linewidth}
\vspace{3mm}
\begin{center}
{\large\bfseries\textcolor{appleblue}{Apple-ID}}
\end{center}
\vspace{2mm}
\begin{itemize}[label=$\square$, itemsep=5pt]
\item Ich verstehe: Apple-ID = \textbf{Schlüssel}
\item Ich weiß, was damit verbunden ist
\item Ich gebe das Passwort \textbf{niemals} weiter
\item Ich weiß, dass Apple den Code \textbf{nicht} kennt
\end{itemize}
\vspace{3mm}
\end{minipage}%
}
\end{center}
\end{minipage}%
\hfill%
\begin{minipage}[t]{0.31\linewidth}
\begin{center}
\fcolorbox{icloudblue}{icloudblue!10}{%
\begin{minipage}{0.92\linewidth}
\vspace{3mm}
\begin{center}
{\large\bfseries\textcolor{icloudblue}{iCloud}}
\end{center}
\vspace{2mm}
\begin{itemize}[label=$\square$, itemsep=5pt]
\item Ich weiß, was iCloud speichert
\item Ich kann iCloud-Einstellungen finden
\item Ich verstehe \textbf{Synchronisation}
\item Ich kenne \textbf{icloud.com}
\end{itemize}
\vspace{3mm}
\end{minipage}%
}
\end{center}
\end{minipage}%
\hfill%
\begin{minipage}[t]{0.31\linewidth}
\begin{center}
\fcolorbox{findmygreen}{findmygreen!10}{%
\begin{minipage}{0.92\linewidth}
\vspace{3mm}
\begin{center}
{\large\bfseries\textcolor{findmygreen}{Weitere Dienste}}
\end{center}
\vspace{2mm}
\begin{itemize}[label=$\square$, itemsep=5pt]
\item Ich weiß, was der \textbf{App Store} ist
\item Ich kann ``\textbf{Wo ist?}'' nutzen
\item Ich kenne \textbf{iMessage} vs. SMS
\item Ich weiß, was \textbf{FaceTime} ist
\end{itemize}
\vspace{3mm}
\end{minipage}%
}
\end{center}
\end{minipage}

\vspace{8mm}

% Gesamtverständnis
\begin{center}
\fcolorbox{black}{lightgray}{%
\begin{minipage}{0.95\linewidth}
\vspace{4mm}
\begin{center}
{\large\bfseries Das Wichtigste}
\end{center}
\vspace{3mm}
\begin{multicols}{2}
\begin{itemize}[label=$\square$, itemsep=6pt]
\item Die \textbf{Apple-ID} verbindet alles
\item Daten sind auf \textbf{allen Geräten} synchron
\item Gekaufte Apps gehören mir auf \textbf{allen Geräten}
\item Ich kann verlorene Geräte \textbf{orten und sperren}
\item Ich verstehe \textbf{iMessage} (blau = kostenlos)
\end{itemize}
\end{multicols}
\vspace{3mm}
\end{minipage}%
}
\end{center}

\vspace{8mm}

% Notfall-Notizen
\begin{center}
\fcolorbox{bordergray}{white}{%
\begin{minipage}{0.95\linewidth}
\vspace{4mm}
\begin{center}
{\large\bfseries Meine Apple-Notizen}\\[3pt]
{\footnotesize (Trage hier nützliche Informationen ein)}
\end{center}
\vspace{4mm}
\renewcommand{\arraystretch}{2}
\begin{tabular}{@{}p{0.45\linewidth}p{0.45\linewidth}@{}}
\textbf{Meine Apple-ID (E-Mail):} & \textbf{iCloud-Speicher:} \\
\dotfill & $\square$ 5 GB ~~ $\square$ 50 GB ~~ $\square$ 200 GB \\[2mm]
\textbf{``Wo ist?'' aktiviert?} & \textbf{iCloud-Fotos aktiviert?} \\
$\square$ Ja ~~~ $\square$ Nein ~~~ $\square$ Weiß nicht & $\square$ Ja ~~~ $\square$ Nein ~~~ $\square$ Weiß nicht \\
\end{tabular}
\renewcommand{\arraystretch}{1}
\vspace{4mm}
\end{minipage}%
}
\end{center}

\vfill

\begin{center}
\footnotesize\textit{Stand: \today{} -- Die Apple-ID ist dein Schlüssel zu allem. Passwort niemals weitergeben!}
\end{center}

\end{document}
